% GENERATED FILE. Edit canonical lessons, then run npm run generate:books.
% canonical-source-hash: fnv1a64:20e5407d39c74bf2
% canonical-lessons: JA-C514-suisoku-suru, JA-C514-yosoku-suru, JA-C514-mikomi, JA-C514-yohou, JA-C514-mitoushi

\chapter{Guesses and Forecasts}
\label{ch:ja-guesses-and-forecasts}
\hlchaptermodality{514}

\begin{chapteropening}
\textbf{By the end of this chapter:} \emph{I can talk about what is likely to happen.}

Complete the last lesson of chapter 514: I can talk about what is likely to happen.
\end{chapteropening}

\section[\texorpdfstring{\ja{すいそくする} (suisoku suru)}{すいそくする (suisoku suru)}]{\ja{すいそくする} (suisoku suru) — to guess; to conjecture}

\label{lesson:JA-C514-suisoku-suru}



\begin{quote}
Before the new one: say the Japanese for to adjust, then the Japanese for arrangements.
\end{quote}

\subsection*{You'll want to know: \ja{すいそくする}}
\textbf{\ja{すいそくする}} — \emph{suisoku suru} — \textquotedblleft{}to guess; to conjecture\textquotedblright{}.

\begin{quote}
\textbf{\ja{もし} \ja{まよったら、すいそくしないで} \ja{ください。}} — \emph{moshi mayottara, suisokushinaide kudasai} — \textquotedblleft{}If you are unsure, please do not guess.\textquotedblright{}
\end{quote}

Polite \textbf{\ja{すいそくします}}, \emph{suisokushimasu}; a \ja{する} verb. Past (polite): \textbf{\ja{せんせいが} \ja{こたえを} \ja{すいそくしました。}}, \emph{sensei ga kotae o suisokushimashita}.

\begin{cousinweb}
\textbf{\ja{すいそくする}} (\emph{suisoku suru}) is the Sino-Japanese (kango) noun \emph{suisoku} with the native verb \emph{suru}, \textquotedblleft{}to do\textquotedblright{}: two Chinese characters read with their on'yomi, readings borrowed from Chinese. They are \emph{sui}, \textquotedblleft{}conjecture\textquotedblright{} (Middle Chinese \emph{tsyhwij}) and \emph{soku}, \textquotedblleft{}fathom\textquotedblright{} (\emph{tsrhik}) (KANJIDIC2; Middle Chinese in Baxter's transcription).
\end{cousinweb}

\begin{grammarlens}[title={What you've built}]
One more word for saying what happens if.
\end{grammarlens}

\subsection*{Your turn}
\begin{itemize}
\raggedright
  \item \emph{Say it:} \emph{suisoku suru}
  \item \emph{Say it:} \emph{suisoku suru}, once more
  \item \emph{Say it:} \emph{moshi mayottara, suisokushinaide kudasai} — If you are unsure, please do not guess.
  \item \emph{Recall:} say the Japanese for to adjust, then the Japanese for arrangements, then say \emph{suisoku suru} again
\end{itemize}

\begin{tcolorbox}[breakable,colback=teal!4,colframe=teal!35!black,title={Before you move on}]
What does \ja{すいそくする} mean? (\textquotedblleft{}To guess; to conjecture\textquotedblright{}.) Say it once more.
\end{tcolorbox}
\section[\texorpdfstring{\ja{よそくする} (yosoku suru)}{よそくする (yosoku suru)}]{\ja{よそくする} (yosoku suru) — to predict; to forecast}

\label{lesson:JA-C514-yosoku-suru}



\begin{quote}
Before the new one: say the Japanese for arrangements, then the Japanese for to guess.
\end{quote}

\subsection*{You'll want to know: \ja{よそくする}}
\textbf{\ja{よそくする}} — \emph{yosoku suru} — \textquotedblleft{}to predict; to forecast\textquotedblright{}.

\begin{quote}
\textbf{\ja{もし} \ja{ゆきを} \ja{よそくしたら、しあいを} \ja{ちゅうしします。}} — \emph{moshi yuki o yosokushitara, shiai o chūshi shimasu} — \textquotedblleft{}If they forecast snow, they will call off the match.\textquotedblright{}
\end{quote}

Polite \textbf{\ja{よそくします}}, \emph{yosokushimasu}; a \ja{する} verb. Past (polite): \textbf{\ja{ちちが} \ja{ゆきを} \ja{よそくしました。}}, \emph{chichi ga yuki o yosokushimashita}.

\begin{cousinweb}
\textbf{\ja{よそくする}} (\emph{yosoku suru}) is the Sino-Japanese (kango) noun \emph{yosoku} with the native verb \emph{suru}, \textquotedblleft{}to do\textquotedblright{}: two Chinese characters read with their on'yomi, readings borrowed from Chinese. They are \emph{yo}, \textquotedblleft{}beforehand\textquotedblright{} (Middle Chinese \emph{yoH}) and \emph{soku}, \textquotedblleft{}fathom\textquotedblright{} (\emph{tsrhik}) (KANJIDIC2; Middle Chinese in Baxter's transcription).
\end{cousinweb}

\begin{grammarlens}[title={What you've built}]
One more word for saying what happens if.
\end{grammarlens}

\subsection*{Your turn}
\begin{itemize}
\raggedright
  \item \emph{Say it:} \emph{yosoku suru}
  \item \emph{Say it:} \emph{yosoku suru}, once more
  \item \emph{Say it:} \emph{moshi yuki o yosokushitara, shiai o chūshi shimasu} — If they forecast snow, they will call off the match.
  \item \emph{Recall:} say the Japanese for arrangements, then the Japanese for to guess, then say \emph{yosoku suru} again
\end{itemize}

\begin{tcolorbox}[breakable,colback=teal!4,colframe=teal!35!black,title={Before you move on}]
What does \ja{よそくする} mean? (\textquotedblleft{}To predict; to forecast\textquotedblright{}.) Say it once more.
\end{tcolorbox}
\section[\texorpdfstring{\ja{みこみ} (mikomi)}{みこみ (mikomi)}]{\ja{みこみ} (mikomi) — a prospect; a hope (of)}

\label{lesson:JA-C514-mikomi}



\begin{quote}
Before the new one: say the Japanese for to guess, then the Japanese for to predict.
\end{quote}

\subsection*{You'll want to know: \ja{みこみ}}
\textbf{\ja{みこみ}} — \emph{mikomi} — \textquotedblleft{}a prospect; a hope (of)\textquotedblright{}.

\begin{quote}
\textbf{\ja{みこみは} \ja{ちりょうに} \ja{よります。}} — \emph{mikomi wa chiryō ni yorimasu} — \textquotedblleft{}The prospects depend on the treatment.\textquotedblright{}
\end{quote}

\begin{cousinweb}
\textbf{\ja{みこみ}} (\emph{mikomi}) is a native Japanese compound (wago), inherited rather than borrowed from Chinese: built from \emph{mi}, \textquotedblleft{}seeing\textquotedblright{} and \emph{komi}, \textquotedblleft{}putting in\textquotedblright{} (JMdict, KANJIDIC2).
\end{cousinweb}

\begin{grammarlens}[title={What you've built}]
One more word for saying what happens if.
\end{grammarlens}

\subsection*{Your turn}
\begin{itemize}
\raggedright
  \item \emph{Say it:} \emph{mikomi}
  \item \emph{Say it:} \emph{mikomi}, once more
  \item \emph{Say it:} \emph{mikomi wa chiryō ni yorimasu} — The prospects depend on the treatment.
  \item \emph{Recall:} say the Japanese for to guess, then the Japanese for to predict, then say \emph{mikomi} again
\end{itemize}

\begin{tcolorbox}[breakable,colback=teal!4,colframe=teal!35!black,title={Before you move on}]
What does \ja{みこみ} mean? (\textquotedblleft{}A prospect; a hope (of)\textquotedblright{}.) Say it once more.
\end{tcolorbox}
\section[\texorpdfstring{\ja{よほう} (yohō)}{よほう (yohō)}]{\ja{よほう} (yohō) — a forecast}

\label{lesson:JA-C514-yohou}



\begin{quote}
Before the new one: say the Japanese for to predict, then the Japanese for a prospect.
\end{quote}

\subsection*{You'll want to know: \ja{よほう}}
\textbf{\ja{よほう}} — \emph{yohō} — \textquotedblleft{}a forecast\textquotedblright{}.

\begin{quote}
\textbf{\ja{もし} \ja{よほうが} \ja{あたったら、あしたは} \ja{ゆきでしょう。}} — \emph{moshi yohō ga atattara, ashita wa yuki deshō} — \textquotedblleft{}If the forecast is right, it will snow tomorrow.\textquotedblright{}
\end{quote}

\begin{cousinweb}
\textbf{\ja{よほう}} (\emph{yohō}) is Sino-Japanese (kango): two Chinese characters read with their on'yomi, readings borrowed from Chinese. They are \emph{yo}, \textquotedblleft{}beforehand\textquotedblright{} (Middle Chinese \emph{yoH}) and \emph{hō}, \textquotedblleft{}report\textquotedblright{} (\emph{pawH}), the same two as at the end of \emph{tenkiyohō}, \textquotedblleft{}the weather forecast\textquotedblright{} (KANJIDIC2; Middle Chinese in Baxter's transcription).
\end{cousinweb}

\begin{grammarlens}[title={What you've built}]
One more word for saying what happens if.
\end{grammarlens}

\subsection*{Your turn}
\begin{itemize}
\raggedright
  \item \emph{Say it:} \emph{yohō}
  \item \emph{Say it:} \emph{yohō}, once more
  \item \emph{Say it:} \emph{moshi yohō ga atattara, ashita wa yuki deshō} — If the forecast is right, it will snow tomorrow.
  \item \emph{Recall:} say the Japanese for to predict, then the Japanese for a prospect, then say \emph{yohō} again
\end{itemize}

\begin{tcolorbox}[breakable,colback=teal!4,colframe=teal!35!black,title={Before you move on}]
What does \ja{よほう} mean? (\textquotedblleft{}A forecast\textquotedblright{}.) Say it once more.
\end{tcolorbox}
\section[\texorpdfstring{\ja{みとおし} (mitōshi)}{みとおし (mitōshi)}]{\ja{みとおし} (mitōshi) — an outlook; a forecast}

\label{lesson:JA-C514-mitoushi}



\begin{quote}
Before the new one: say the Japanese for a prospect, then the Japanese for a forecast.
\end{quote}

\subsection*{You'll want to know: \ja{みとおし}}
\textbf{\ja{みとおし}} — \emph{mitōshi} — \textquotedblleft{}an outlook; a forecast\textquotedblright{}.

\begin{quote}
\textbf{\ja{みとおしは} \ja{てんきに} \ja{よります。}} — \emph{mitōshi wa tenki ni yorimasu} — \textquotedblleft{}The outlook depends on the weather.\textquotedblright{}
\end{quote}

\begin{cousinweb}
\textbf{\ja{みとおし}} (\emph{mitōshi}) is a native Japanese compound (wago), inherited rather than borrowed from Chinese: built from \emph{mi}, \textquotedblleft{}seeing\textquotedblright{} and \emph{tōshi}, \textquotedblleft{}passing through\textquotedblright{} (JMdict, KANJIDIC2).
\end{cousinweb}

\begin{grammarlens}[title={What you've built}]
One more word for saying what happens if.
\end{grammarlens}

\subsection*{Your turn}
\begin{itemize}
\raggedright
  \item \emph{Say it:} \emph{mitōshi}
  \item \emph{Say it:} \emph{mitōshi}, once more
  \item \emph{Say it:} \emph{mitōshi wa tenki ni yorimasu} — The outlook depends on the weather.
  \item \emph{Recall:} say the Japanese for a prospect, then the Japanese for a forecast, then say \emph{mitōshi} again
\end{itemize}

\begin{tcolorbox}[breakable,colback=teal!4,colframe=teal!35!black,title={Before you move on}]
What does \ja{みとおし} mean? (\textquotedblleft{}An outlook; a forecast\textquotedblright{}.) Say it once more.
\end{tcolorbox}
